\documentclass[11pt,a4paper]{article}

% Typography & Layout.
\usepackage[utf8]{inputenc}
\usepackage[T1]{fontenc}
\usepackage{mathpazo}
\usepackage[margin=1.25in]{geometry}
\usepackage{titlesec}
\usepackage{setspace}
\usepackage{microtype}

% Formatting.
\titleformat{\section}{\Large\bfseries}{\thesection}{1em}{}
\titleformat{\subsection}{\large\bfseries}{\thesubsection}{1em}{}
\onehalfspacing

% Document...
\begin{document}

\begin{center}
    \vspace*{2em}
    {\Huge \textbf{Beyond Coexistence}}\\[0.5em]
    {\Large \textit{The Mechanics of a Plural Society}}\\[2em]

    {\large \textbf{Juan S. A. Siagian}}\\[0.5em]
    {\normalsize Sunken Court, ITB \hspace{0.5em} $\cdot$ \hspace{0.5em} \today}\\[3em]
\end{center}

\section*{The Inherent Complexity}
Humans are fundamentally shaped by their lived experiences. And because no two lives are identical, no two minds are perfectly aligned. But the true hurdle of a plural society is not diversity itself, but the staggering complexity that diversity generates. We each possess distinct thoughts, needs, and preferences, making it incredibly difficult to imagine a social architecture capable of accommodating this scale of human variance.

\end{document}